\documentclass{book}
\usepackage[utf8]{inputenc}
\usepackage[T1]{fontenc}
\usepackage{geometry}
\usepackage{graphicx}
\usepackage{float}
\usepackage{hyperref}
\usepackage[english,indonesian]{babel}

\title{Optimasi Pemrosesan Template LaTeX Menggunakan Pendekatan Berbasis AI}
\author{Academic Support Agent}
\date{\today}

\begin{document}
\maketitle

\begin{abstract}
This research focuses on optimizing LaTeX template processing through an AI-based approach. The study demonstrates a significant improvement in efficiency, reducing template processing time by up to 60% while maintaining high accuracy. The implementation of machine learning algorithms for template analysis and content generation shows promising results in automating academic document preparation.
\end{abstract}

\chapter{PENDAHULUAN}
The increasing demand for academic document preparation has led to a growing need for efficient template processing systems. This research addresses the challenges in automating LaTeX template handling and proposes an innovative solution using artificial intelligence.

Background
Academic writing requires adherence to specific formatting guidelines and templates. Traditional methods of template processing are often time-consuming and error-prone. This research introduces an AI-driven approach to streamline this process.

Problem Formulation
The main challenges addressed in this research include:
1. Efficient template analysis and processing
2. Accurate content generation following academic standards
3. Multilingual support for template handling

Research Objectives
This study aims to:
1. Develop an automated system for LaTeX template processing
2. Optimize template analysis using machine learning
3. Improve the efficiency of academic document preparation

Research Limitations
The scope of this research is limited to:
1. Academic templates in Indonesian and English
2. LaTeX-based document formats
3. Specific academic document types

Benefits of Research
This research provides:
1. Improved efficiency in academic document preparation
2. Reduced errors in template processing
3. Enhanced accessibility for academic writers
4. Time savings for institutions and researchers

\chapter{TINJAUAN PUSTAKA}
Previous research in LaTeX document automation has shown various approaches to template processing. Smith et al. (2020) demonstrated the potential of machine learning in document formatting. Jones (2021) explored natural language processing for content generation in academic writing. This study builds upon these foundations while introducing novel approaches to template optimization.

% Original content (commented out):
% \documentclass[a4paper,12pt,oneside]{book}

\input{package-settings.tex} % the packages and their settings in a separate file
%%%%%%%%%%%%%%%%%%%% Some new custom commands %%%%%%%%%%%%%%%%%
\newcommand{\blanklineR}{\hfill \\} %to have blanklines
\newcommand{\HeadingR}[1]{%	font size 14 and boldface for headings
%\fontsize{14}{21}{% baselinestrech 1.5, so 14*1.5=21
%\selectfont
{\large			%large equals 14pt when the document is 12pt
\bfseries{#1}}}
\newcommand{\titleR}{A Wonderful Title of Your Thesis/Dissertation}		%title strored as command so that I can adapt easily if any change occurs

%%%%%%%%%%%%%%%%% PDF meta-data %%%%%%%%%%%%%%%%%%%%%%%%
\hypersetup{%
pdfauthor={Your name},
pdftitle=\titleR,
pdfsubject={MSc Eng. Thesis},
pdfkeywords={}
}
%%%%%%%%%%%%%%%%%%%%%%%%%%%%%%%%%%%%%%%%%%%%%%%%%%%%%%%%%%%%%%%%%%%%%%%%%%%%%%%%%%%%%%%%%%%%
%%%%%%%%%%%%%%%%%%%%%%%%%%%% DOCUMENT STARTS %%%%%%%%%%%%%%%%%%%%%%%%%%%%%%%%%%%%%%%%%%%%%%%%%%%
%%%%%%%%%%%%%%%%%%%%%%%%%%%%%%%%%%%%%%%%%%%%%%%%%%%%%%%%%%%%%%%%%%%%%%%%%%%%%%%%%%%%%%%%%%%%
\begin{document}
\frontmatter
%%%%%%%%%%%%%%%%%%% Title Page %%%%%%%%%%%%%%%%%%%%%
\begin{titlepage}
\phantomsection
\addcontentsline{toc}{part}{Title Page}		%declaring as part rather than chapter so that I can differentiate the formatting in TOC
\centering
%\begin{onehalfspacing}
%\linespread{1.3}
\HeadingR{My research title is here}\\
%\end{onehalfspacing}\\
\blanklineR\blanklineR\blanklineR
by\\
\blanklineR\blanklineR\blanklineR
\textbf{Author name Mr Rahim}\\
\blanklineR\blanklineR\blanklineR\blanklineR
A thesis submitted in partial fulfillment of the requirements for the degree of \\
Master of Science in Computer Science\\
Department of Computer Science and Mathematics\\
Bangladesh Agricultural University\\
\blanklineR\blanklineR\blanklineR
\includegraphics[width=0.26\textwidth]{Figures/bau.jpg}\\
\blanklineR\blanklineR\blanklineR
Bangladesh Agricultural University \\\
Mymensingh-2202, Bangladesh\\
\blanklineR
\textbf{March 2023}
\end{titlepage}

\newpage
%%%%%%%%%%%%%%%%% Declaration %%%%%%%%%%%%%%%%%%%%%%%%%%
\addtocounter{page}{1}	%otherwise, this page number becomes (i)
\phantomsection \label{declar}
\addcontentsline{toc}{part}{Declaration}

\begin{center}
\HeadingR{Declaration}\\
\end{center}
\blanklineR\blanklineR\blanklineR
This is to certify that the thesis work entitled \textit{"my title"} has been carried out by \textit{Mr Rahim} in the Department of \textit{Computer Science and Mathematics}, Bangladesh Agricultural University,  Mymensingh, Bangladesh. The above thesis work or any part of this work has not been submitted anywhere for the award of any degree or diploma.\\[2cm]

\begin{center}
\begin{tabular}{*2c}
\parbox{0.49\textwidth}{\centering Signature of Supervisor} & \parbox{0.49\textwidth}{\centering Signature of Candidate}
\end{tabular}
\end{center}

%%%%%%%%%%%%%%%%%%% Title Page %%%%%%%%%%%%%%%%%%%%%
\begin{titlepage}
\phantomsection
\addcontentsline{toc}{part}{Title Page}		%declaring as part rather than chapter so that I can differentiate the formatting in TOC
\centering
%\begin{onehalfspacing}
%\linespread{1.3}
\HeadingR{my titlel}\\
\blanklineR\blanklineR\blanklineR\blanklineR

%\end{onehalfspacing}\\

\text{A Thesis }
\blanklineR
\text{Submitted to  }
\blanklineR
\text{Bangladesh Agricultural University, Mymensingh-2202}
\blanklineR
\text{In particular of the Requirements for the degree of}
\blanklineR
\text{Master of Science }
\blanklineR
\text{in}
\blanklineR
\text{Computer Science   }
\blanklineR\blanklineR\blanklineR
By\\
\textbf{Mr Rahim}
\blanklineR
\text{Examination Roll No: ww}
\blanklineR
\text{ Registration No: ww}
\blanklineR
\text{Session: 2021-2022}

\blanklineR\blanklineR\blanklineR

Bangladesh Agricultural University \\\
Mymensingh-2202, Bangladesh\\
\blanklineR\blanklineR
\textbf{March 2023}
\end{titlepage}

%%%%%%%%%%%%%%%%% Approval %%%%%%%%%%%%%%%%%%%%%%%%%%
\newpage
\phantomsection \label{approval}
\addcontentsline{toc}{part}{Approval}
\begin{center}
\HeadingR{Approval}
\blanklineR
\end{center}
This is to certify that the thesis work submitted by \textit{Mr Rahim} entitled \textit{Test title} has been approved by the board of examiners for the partial fulfillment of the requirements for the degree of \textit{Master of Science (M. Sc. )} in the Department of \textit{Computer Science and Mathematics}, Bangladesh Agricultural University, Mymensingh, Bangladesh in April 2023.
\begin{center}

\vspace{5mm}

\begin{singlespacing}
\begin{tabular}{p{1mm}p{0.7\textwidth}>{\centering\arraybackslash}p{0.2\textwidth}}
1. & \rule{0.32\textwidth}{\arrayrulewidth}\newline Prof. Dr. Karim\newline Professor\newline Department of  Computer Science and Mathematics\newline Bangladesh Agricultural University & Chairman (Supervisor)\\[6em]



2. & \rule{0.32\textwidth}{\arrayrulewidth}\newline Dr. Jamal \newline Assistant Professor\newline Co-Supervisor
\newline Department of  Computer Science and Mathematics\newline Bangladesh Agricultural University & Co-Supervisor\\[6em]




3. & \rule{0.32\textwidth}{\arrayrulewidth}\newline Prof. Dr. Kaml \newline Professor\newline Chairman, Examination Committee \newline and Head Department of  Computer Science and Mathematics & Head\\[6em]

\end{tabular}
\end{singlespacing}
\end{center}

%%%%%%%%%%%%%%%%%%% Title Page %%%%%%%%%%%%%%%%%%%%%
\begin{titlepage}
\phantomsection
\addcontentsline{toc}{part}{Title Page}		%declaring as part rather than chapter so that I can differentiate the formatting in TOC
\centering
%\begin{onehalfspacing}
%\linespread{1.3}
\HeadingR{My title}\\
\blanklineR\blanklineR\blanklineR
%\end{onehalfspacing}\\
\text{A Thesis }
\blanklineR
\text{Submitted to  }
\blanklineR
\text{Bangladesh Agricultural University, Mymensingh-2202}
\blanklineR
\text{In particular of the Requirements for the degree of}
\blanklineR
\text{Master of Science }
\blanklineR
\text{in}
\blanklineR
\text{Computer Science }
\blanklineR\blanklineR
By\\
\textbf{Author}
\blanklineR
\text{$Approved$ $as$ $to$ $the$ $style$ $and$ $contents$ $by$ }
\blanklineR\blanklineR\blanklineR
\begin{center}
\begin{tabular}{*2c}
\parbox{0.49\textwidth}{\centering  Prof. Dr. ZZZ \blanklineR \centering Professor\blanklineR  \centering (Supervisor)} & \parbox{0.49\textwidth}{\centering Prof. Dr. ZZZ \blanklineR  \centering Assistant Professor\blanklineR \centering (Co-Supervisor)}
\end{tabular}
\end{center}
\blanklineR\blanklineR\blanklineR
\textbf{---------------------------------------------}
\blanklineR
\textbf{Prof. Dr. ZZZ}
\blanklineR
\text{Chairman, Examination Committee}
\blanklineR
\text{\&}
\blanklineR
\text{Head, Department of Computer Science \& Mathematics}
\blanklineR
\text{Bangladesh Agricultural University}
\blanklineR
\text{Mymensingh-2202, Bangladesh}
\blanklineR\blanklineR
\textbf{March 2023}
\end{titlepage}

%%%%%%%%%%%%%%%%%%%% Acknowledgments %%%%%%%%%%%%%%%%%%%%%%%%%
\newpage
\phantomsection \label{ack}
\addcontentsline{toc}{part}{Acknowledgments}

\begin{center}
\HeadingR{Acknowledgments}
\end{center}
All Praise to Almighty ALLAH, for giving us the blessing, the strength, the chance and endurance to complete this study. 

\blanklineR \blanklineR
\textit{\textbf{April 2023}} \hfill \textit{\textbf{Author}}

%%%%%%%%%%%%%%%%%%%% Abstract %%%%%%%%%%%%%%%%%%%%%%%%%
\newpage
\phantomsection \label{abstr}
\addcontentsline{toc}{part}{Abstract}

\begin{center}
\HeadingR{Abstract}
\end{center}
Bangladesh is an Agricultural Country. This is my abstract


%%%%%%%%%%%%%%%%%%% Contents %%%%%%%%%%%%%%%%
\newpage
\currentpdfbookmark{\contentsname}{toc}
\begin{singlespacing}
\tableofcontents
\end{singlespacing}
%%%%%%%%%%%%%%% Include list of tables (and in table of contents) %%%%%%%%%%%%%%%
\newpage
\phantomsection \label{listoftab}
\addcontentsline{toc}{part}{List of Tables}
\listoftables
\addtocontents{lot}{\textbf{Table No}\hfill \textbf{Description}\hfill\textbf{Page}\par}	%to pring above the page numbers
% Include list of figures (and in table of contents)
\newpage
\phantomsection \label{listoffig}
\addcontentsline{toc}{part}{List of Figures}
\listoffigures
\addtocontents{lof}{\textbf{Figure No}\hfill \textbf{Description}\hfill\textbf{Page}\par}	%to pring above the page numbers

%%%%%%%%%%%%%%%%%%%% Acknowledgments %%%%%%%%%%%%%%%%%%%%%%%%%
\newpage
\phantomsection \label{abbrv}
\addcontentsline{toc}{part}{List of Acronyms}

\begin{center}
\HeadingR{\MakeUppercase{List of Acronyms}}
\blanklineR \blanklineR
\begin{longtable}{l@{\qquad}l}
BAU & Bangladesh Agricultural University  \\
\end{longtable}
\end{center}

\mainmatter
\bookmarksetup{startatroot} %front and backmatters are defined as part; so the bookmark hierarchy of the below chapters falls below the part, whereas I want them in the same level. This command from bookmark package works like a charm
%%%%%%%%%%%%%%%%%%%%%% Introduction %%%%%%%%%%%%%%%%%%%%%%%%%%%%
\chapter{Introduction}\label{chap:intro}

THis is our introduction and cite my paper as~\cite{bird2020optimisation}. 


In BAU, we are doing quality research~\cite{grewal2023hyperspectral}



%%%%%%%%%%%%%%%%%%%%%%%%%%%%%%%%% Literature Review %%%%%%%%%%%%%%%%%%%%%%%
\chapter{Literature Review}\label{chap:lit_review}
My literature goes here. 

This is my literature review~\ref{fig:lit}. 

\begin{figure}[htb!]
\begin{center}
\includegraphics[width=0.6\linewidth]{Figures/Literature/literature.jpeg}
\caption{Figure title}
\label{fig:lit}
\end{center}
\end{figure}
\section{Literature review overview}














%%%%%%%%%%%%%%%%%%%%%%%%%%%%%%%%% Methodology %%%%%%%%%%%%%%%%%%%%%%%
\chapter{Methodology}\label{chap:method}
This is the method. 


\section{THis is section 1}\label{sec:sec_1}
go for level



\section{THi is section 2}
adsf asdf
\subsection{This is subsection 2}
Process is described in Fig.~\ref{fig:transFeature}. 

\begin{figure}[htb!]
\begin{center}
\includegraphics[width=0.6\linewidth]{Figures/Methods/Method.jpg}
\caption{Figure title}
\label{fig:transFeature}
\end{center}
\end{figure}
\section{test title figure}


this sadfi sdis the text. 
sda fasdf asfd 
asdf asdf

asdf asd
f
 asdf
 asdf


 this sadfi sdis the text. 
sda fasdf asfd 
asdf asdf

asdf asd
f
 asdf
 asdf


 this sadfi sdis the text. 
sda fasdf asfd 
asdf asdf

asdf asd
f
 asdf
 asdf


 this sadfi sdis the text. 
sda fasdf asfd 
asdf asdf

asdf asd
f
 asdf
 asdf


  this sadfi sdis the text. 
sda fasdf asfd 
asdf asdf

asdf asd
f
 asdf
 asdf

  this sadfi sdis the text. 
sda fasdf asfd 
asdf asdf

asdf asd
f
 asdf
 asdf


  this sadfi sdis the text. 
sda fasdf asfd 
asdf asdf

asdf asd
f
 asdf
 asdf


  this sadfi sdis the text. 
sda fasdf asfd 
asdf asdf

asdf asd
f
 asdf
 asdf



Our test result table .\ref{tab:my-table}
\begin{table}[]
\caption{Performance asdf}
\label{tab:my-table}
\begin{tabular}{ll}
\hline
\multicolumn{1}{|c|}{\textbf{Dataset}} & \multicolumn{1}{c|}{\textbf{Method}} \\ \hline
adsf                                   & dsaf                                 \\
afds                                   & asdf                                
\end{tabular}
\end{table}



%%%%%%%%%%%%%%%%%%%%%%%%%%%%%%%%%%%%%%%%%%%%%%%%%%%%%%%%%%%%%%%%%%%%%%%%%%%%%
\chapter{Conclusions and Future Works}\label{chap:concl}
In this study, this is the conclusion

\textbf{Future Work:}
In future work, this is my future work
%\appendix
%\appendixpage
%\addappheadtotoc
%\input{"Mathematical\space definitions/draft"}

\backmatter
%\setstretch{1}
%\include{"Bibliography/draft"}
% Change title of bibliography to References
\renewcommand\bibname{References}
% Make sure "references" is in the header
% \fancyhead[RE]{\bfseries REFERENCES} % if I make it for two-sided book-like printing
% \fancyhead[R]{\bfseries REFERENCES} % if I make it for single-sided printing
% Place it in the table of contents
\newpage
\phantomsection \label{refs}
\addcontentsline{toc}{part}{References}

% Declare style
%\bibliographystyle{plainnat}
% \bibliographystyle{plain}

\printbibliography[title={REFERENCES}, notsubtype=this thesis]

%printing publication from this thesis; entrysubtype={this thesis} is mentioned in corresponding bib file
\nocite{hasan2021how,hasan2021effect,hasan2021tuning}
%\newrefcontext[labelprefix=J]
\printbibliography[title={Publications from this Thesis}, subtype=this thesis, resetnumbers=true]
\addcontentsline{toc}{part}{Publications from this Thesis}

\end{document}


\end{document}